%! Conic Sections — Unit 4, Lesson 4.6 — Homework (Answer Key)
\documentclass[10pt]{article}
\usepackage{precalculus-article}
\usepackage{precalculus-key}   % pulls in -boxes; do NOT also load -boxes
\newcommand{\TallMath}[1]{$\displaystyle #1\rule[-1.4em]{0pt}{3.2em}$}

\begin{document}

\pageheader{Unit 4, Lesson 4.6}{Homework --- Answer Key}
\namedateperiod

\vspace{0.10in}

\begin{notesbox}{Conic Sections Practice --- Key}

\textbf{1.}\quad Identify the conic type, center or vertex, and key measurements.

\vspace{4pt}
(a)\quad $y + 1 = \tfrac{1}{2}(x - 3)^2$

\vspace{4pt}
Type: \ans{parabola} \quad Vertex: \ans{$(3,\,-1)$} \quad $a =$ \ans{$\tfrac{1}{2}$} \quad Opens: \ans{upward}

\vspace{10pt}
(b)\quad $\dfrac{(x+2)^2}{36} + \dfrac{(y-4)^2}{16} = 1$

\vspace{4pt}
Type: \ans{ellipse} \quad Center: \ans{$(-2,\,4)$} \quad $a =$ \ans{$6$} \quad $b =$ \ans{$4$}

\vspace{10pt}
(c)\quad $\dfrac{(x-5)^2}{4} - \dfrac{(y+3)^2}{25} = 1$

\vspace{4pt}
Type: \ans{hyperbola} \quad Center: \ans{$(5,\,-3)$} \quad $a =$ \ans{$2$} \quad $b =$ \ans{$5$}

\vspace{10pt}
\textbf{2.}\quad Write the standard form equation.

\vspace{6pt}
(a)\quad Parabola, vertex $(0,4)$, opening downward, through $(2,0)$.

\vspace{4pt}
\ans{$y - 4 = a(x-0)^2$. Substitute $(2,0)$: $-4 = 4a \Rightarrow a=-1$.
Equation: $y - 4 = -(x)^2$.}

\vspace{10pt}
(b)\quad Parabola, vertex $(3,-1)$, opening left, through $(1,3)$.

\vspace{4pt}
\ans{$x - 3 = a(y+1)^2$. Substitute $(1,3)$: $-2 = 16a \Rightarrow a = -\tfrac{1}{8}$.
Equation: $x - 3 = -\tfrac{1}{8}(y+1)^2$.}

\vspace{10pt}
(c)\quad Ellipse centered at $(-1,2)$, horizontal radius 6, vertical radius 4.

\vspace{4pt}
\ans{$\dfrac{(x+1)^2}{36} + \dfrac{(y-2)^2}{16} = 1$}

\vspace{10pt}
\textbf{3.}\quad Hyperbola centered at $(4,-3)$, $a=5$, $b=2$, opening up and down.

\vspace{6pt}
\ans{$\dfrac{(y+3)^2}{25} - \dfrac{(x-4)^2}{4} = 1$}

\vspace{4pt}
\ans{Asymptotes: $y + 3 = \pm\dfrac{5}{2}(x-4)$}

\vspace{10pt}
\textbf{4.}\quad $9x^2 - 4y^2 = 36$.

\vspace{4pt}
(a)\quad \ans{Divide by 36: $\dfrac{x^2}{4} - \dfrac{y^2}{9} = 1$.}

\vspace{8pt}
(b)\quad \ans{Center: $(0,0)$; $a=2$, $b=3$; opens left and right.
Asymptotes: $y = \pm\dfrac{3}{2}x$.}

\end{notesbox}

\vspace{0.08in}

\begin{notesbox}{AP-Style Practice --- Key}

\textbf{5.}\quad Correct answer: \ans{(B)}

$\dfrac{(x-2)^2}{9} - \dfrac{(y+1)^2}{16} = 1$

\begin{teachernote}
  (A) is an ellipse (uses $+$, not $-$). (B) is correct: left/right hyperbola has positive
  $x$-term; center $(2,-1)$ gives $(x-2)$ and $(y+1)$; $a^2=9$, $b^2=16$. (C) opens up/down
  (positive $y$-term). (D) has wrong center signs: $h=2$ requires $(x-2)$, not $(x+2)$.
\end{teachernote}

\end{notesbox}

\end{document}
